\chapter{مقدمه و نقشه راه}
\label{ch:introduction}

\section{این راهنما برای چه کسی است؟}

اگر قصد تحصیل در خارج از کشور را دارید، چه در مقطع کارشناسی ارشد
یا دکتری این راهنما مسیر کلی را مرحله‌به‌مرحله توضیح می‌دهد.
فرض بر این است که با مفاهیم پایه آشنا هستید اما به یک منبع ساختاریافته
و به‌روز نیاز دارید.
تلاش شده که ضمن حفظ جامعیت، این راهنما در بخش‌های مختلف به منابع بیرونی که توضیحات کامل‌تری درباره مراحل مختلف 
می‌دهند ارجاع داشته باشد تا هم جلوی طولانی شدن راهنما گرفته شود و هم افراد بسته به زمانی که دارند از 
میزان‌های مختلف راهنمایی بهره ببرند.
\section{نقشه کلی مسیر}
مسیر مهاجرت تحصیلی یک مسیر طولانی است که می‌تواند حداقل یک‌سال یا بیشتر طول بکشد. مراحل مختلف در این مسیر
بسته به جنسیت، مقصد، دانشگاه و فاکتور‌های دیگر بسیار متغیر است ولی یک توضیف تسبتا عمومی از مراحل میتواند شامل موارد زیر باشد. این مورد‌ها 
این مورد‌ها دارای ترتیب هستند ولی گاهی مجبور هستید بعضی از آن‌ها را (مثلا زبان و درخواست پذیرش) را به صورت موازی پیش‌ببرید.

\begin{enumerate}
  \item \textbf{تحقیق و مشخص کردن اولویت‌ها}: در این مرحله باید تلاش کنید که از بین گزینه‌های کلی موجود بین قاره‌ها و کشور‌های مختلف
  بر اساس شرایط مالی، درسی و سایر چیزهایی که ممکن است برای شما مهم باشد تعدادی گزینه انتخاب کنید. در این مرحله 
  تلاش نکنید که گزینه نهایی خود را مشخص کنید چون عوامل بی‌شماری در مراحل آینده می‌تواند روی انتخاب شما تاثیر بگذارد. به عنوان مثال ممکن است شما 
  در ابتدا کانادا را به عنوان گزینه اصلی در نظر بگیرید و در ادامه به دلیل تغییر قوانین مهاجرتی مجبور به تغییر گزینه شوید. پس هدف ما در
  این مرحله مشخص کردن انتخاب‌های ممکن و معقول به صورت کلی است و نه انجام انتخاب نهایی.
  \item \textbf{آمادگی مدارک}: یکی از کارهایی که به طور کلی باید قبل از ارسال درخواست‌های پذیرش انجام شود، تهیه یک چک‌لیست از مدارک مختلفی
  است که به آن‌ها نیاز خواهیم داشت. هر دانشگاه برای درخواست پذیرش لیستی از مدارک ارائه می‌دهد که ممکن است در جزئیات با سایر دانشگاه‌ها فرق کند، 
  مثلا شاید یک دانشگاه دانشجویان را ملزم به استفاده از فرمت خاصی برای ارائه \lr{CV} نماید. اما به طور کلی مدارکی که تقریبا همه دانشگاه‌ها 
  در هنگام درخواست پذیرش از شما خواهند خواست شامل این موارد است: ریزنمرات، \lr{CV}، انگیزه‌نامه (\lr{Statement of Purpose})و توصیه‌نامه.
  دانشگاه‌ها بر اساس همین مدارک اقدام به سنجیدن شما نسبت به سایر \lr{Applicant} ها می‌کنند بنابرین کیفیت و چیزهایی که در این مدارک ارائه 
  می‌دهید بسیار مهم است..
  \item \textbf{آزمون‌ها}: برای مقایسه راحت و رسیدن به یک معیار عددی در ارزیابی دانشجویان در بعضی زمینه‌ها \textbf{سفارت و دانشگاه} از 
  شما مدارکی جهت اثبات توانایی در زمینه‌های مختلف می خواهند. پر تکرار ترین این مدارک نتیجه آزمون زبان است. آزمون‌های \textbf{TOEFL} و \textbf{IELTS}
  معتبرترین آزمون‌های زبان انگلیسی هستند و اکثر قریب به اتفاق دانشگاه‌های ارائه دهنده پروگرم‌های انگلیسی زبان همین دو آزمون را به عنوان اثبات توانایی زبان قبول دارند. 
  آمادگی برای این آزمون‌ها در صورت نداشتن پیشینه قوی در زبان انگلیسی می‌تواند زمانبر باشد و باید برای آن ها از مدتی قبل بسته به توانایی خود برنامه‌ریزی کنید.
  \item \textbf{مکاتبه و درخواست پذیرش}: بسته به مقطع و کشور هدف شما، مراحل مختلفی برای درخواست پذیرش وجود دارد. به طور کلی در این مرحله ممکن است به مکاتبه با 
  اساتید جهت یافتن استاد راهنما/سرپرست، ثبت مدارک در سایت دانشگاه، تهیه و ارائه پروپوزال و مواردی از این دست داشته باشید. بنابرین پس از مشخص کردن کشور‌ها و دانشگاه‌های 
  مقصد خود باید زمان‌بندی و ددلاین‌های انجام هر یک از مراحل پذیرش آن‌ها را یادداشت و بر اساس آن ها مدارک خود را آماده کرده و به انجام آن مراحل بپردازید. 0
  \item \textbf{ویزا}: پس از مشخص شدن نتایج پذیرش‌های شما و انتخاب کردن گزینه‌ یا گزینه‌های نهایی خود در این مرحله شما از دولت کشور 
  مقصد درخواست اجاره برای اقامت طولانی مدت با هدف تحصیل در یک مقطع خاص را خواهید داشت. این مرحله نیاز به ارائه مدارک خاصی با توجه 
  به قوانین هر کشور دارد و بعضا قوانین و مدارک هر سال به روز می‌شود.
  \item \textbf{آمادگی قبل از سفر}: بیمه، مسکن، بانک، بلیط
  \item \textbf{ورود و استقرار}: فرودگاه، ثبت‌نام دانشگاه، \lr{SIN/SSN} و...
\end{enumerate}

\begin{infobox}[title={ساختار راهنما}]
  هر فصل یکی از این مراحل را پوشش می‌دهد. در پایان، فصل
  \cref{ch:resources} لیست جامعی از پیوندهای مفید دارد.
\end{infobox}

\section{اصطلاحات رایج}

\begin{table}[htbp]
  \centering
  \caption{واژه‌های پرکاربرد}
  \label{tab:glossary}
  \begin{tabular}{@{}p{0.35\textwidth}p{0.55\textwidth}@{}}
    \toprule
    \textbf{اصطلاح} & \textbf{معنی} \\
    \midrule
    \lr{Conditional Offer} & پذیرش مشروط (مثلاً تا ارائه مدرک زبان) \\
    \lr{CAS / I-20} & سند پذیرش رسمی برای درخواست ویزا \\
    \lr{Proof of Funds} & اثبات تمکن مالی \\
    \lr{Biometrics} & اثر انگشت و عکس برای ویزا \\
    \bottomrule
  \end{tabular}
\end{table}

\section{قبل از شروع}

\begin{checklistbox}
  \begin{checklist}
    \item کشور(های) هدف را مشخص کنید
    \item بودجه تقریبی (شهریه + زندگی) را برآورد کنید
    \item مهلت‌های \lr{Deadline} دانشگاه‌ها را در تقویم ثبت کنید
    \item یک پوشه (فیزیکی یا دیجیتال) برای مدارک بسازید
  \end{checklist}
\end{checklistbox}

\begin{tipbox}
  برای هر کشور یک زیرپوشه در \lr{Google Drive} یا \lr{Notion} بسازید
  و لینک‌های رسمی را همان‌جا نگه دارید. بعداً برای تبدیل به وبلاگ
  هم از همین ساختار استفاده می‌کنید.
\end{tipbox}
